% Originally: new file (chapter 18 had no solutions section, having had no exercises).

\section{Solutions}
\label{sec:geodesics_solutions}

\begin{exercisesolution}[Solution to \cref{ex:ai_fundamental_lemma}]
% Generator: Claude Opus 5 (High)
Suppose $h \not\equiv 0$, so that $h(t_0) \neq 0$ for some $t_0$. Since $h$ is continuous, the
set where $h \neq 0$ is open in $[0,1]$, so we may take $t_0 \in (0,1)$: if $h$ were non-zero
only at an endpoint it would still be non-zero just inside.

Choose a component $i$ with $h_i(t_0) \neq 0$. Replacing $h$ by $-h$ if necessary, which changes
neither the hypothesis nor the conclusion, assume $h_i(t_0) =: 2m > 0$. By continuity there is a
$\delta > 0$ with $(t_0-\delta,\,t_0+\delta) \subseteq (0,1)$ and
\[h_i(t) > m \qquad\text{for all } t \in (t_0-\delta,\,t_0+\delta).\]
Now take the bump function of \cref{def:smooth_bump_function} on that interval,
\[\varphi(t) := \begin{cases}
    \exp\!\left(\dfrac{-1}{\delta^2 - (t-t_0)^2}\right), & \lvert t-t_0\rvert < \delta, \\[2ex]
    0, & \text{otherwise,}
  \end{cases}\]
which is smooth on $\mathbb{R}$, non-negative, strictly positive on $(t_0-\delta,t_0+\delta)$ and
zero outside it, and set $g := \varphi\,e_i$. Then $g$ is smooth with $g(0) = g(1) = 0$, so it is
an admissible test function, and
\[\int_0^1 \big\langle h(t),\,g(t)\big\rangle\,dt
  = \int_{t_0-\delta}^{t_0+\delta} h_i(t)\,\varphi(t)\,dt
  \ \geq\ m\int_{t_0-\delta}^{t_0+\delta}\varphi(t)\,dt \ > \ 0,\]
the last inequality because $\varphi$ is continuous, non-negative and not identically zero on the
interval. This contradicts the hypothesis. Consequently $h \equiv 0$.

\begin{remark}[Where each hypothesis is spent]
\label{rem:fundamental_lemma_hypotheses}
Two ingredients carry the proof, and it is worth seeing which does what. \emph{Continuity} of $h$
turns a single non-zero value into a whole interval of non-zero values of one sign, and without
it the statement is false: a function equal to $1$ at one point and $0$ elsewhere integrates to
zero against everything. \emph{Smoothness} of the bump is what keeps $g$ admissible, and it is
the only reason the test class can be shrunk from all continuous $g$ to all smooth $g$ without
weakening the conclusion. The conclusion itself is pointwise, which is the whole content: the
hypothesis says $h$ vanishes on average against every test function, and the lemma upgrades that
to vanishing everywhere.
\end{remark}
\end{exercisesolution}

\begin{exercisesolution}[Solution to \cref{ex:ai_variation_without_constant_speed}]
% Generator: Claude Opus 5 (High)
\begin{enumerate}[label=\textbf{(\alph*)}]
  \item By the second identity of \cref{lem:inner_product_derivative_identities} applied to
    $\alpha := c_\varepsilon'$, and differentiating under the integral sign,
    \[\left.\frac{d}{d\varepsilon}\right|_{\varepsilon=0} L(c_\varepsilon)
      = \int_0^1 \frac{\big\langle c'(t),\ \partial_\varepsilon\big|_{\varepsilon=0}
        c_\varepsilon'(t)\big\rangle}{\lvert c'(t)\rvert}\,dt
      = \int_0^1 \left\langle \frac{c'(t)}{\lvert c'(t)\rvert},\ g'(t)\right\rangle dt.\]
    Nothing was pulled out of the integral, which is the one difference from the proof of
    \cref{prop:geodesic_straight_line}. Since $c \in C^2$ and $\lvert c'\rvert > 0$, the unit
    tangent $c'/\lvert c'\rvert$ is $C^1$, so integration by parts is available:
    \[\left.\frac{d}{d\varepsilon}\right|_{\varepsilon=0} L(c_\varepsilon)
      = \left[\left\langle \frac{c'}{\lvert c'\rvert},\,g\right\rangle\right]_0^1
        - \int_0^1 \left\langle \left(\frac{c'}{\lvert c'\rvert}\right)'(t),\ g(t)\right\rangle dt
      = -\int_0^1 \left\langle \left(\frac{c'}{\lvert c'\rvert}\right)'(t),\ g(t)\right\rangle dt,\]
    the boundary term vanishing because $g(0) = g(1) = 0$.
  \item The left-hand side is zero for every admissible $g$, and
    $\big(c'/\lvert c'\rvert\big)'$ is continuous, so \cref{ex:ai_fundamental_lemma} applies with
    $h := \big(c'/\lvert c'\rvert\big)'$ and gives $\big(c'/\lvert c'\rvert\big)' \equiv 0$.
  \item Therefore $c'/\lvert c'\rvert \equiv u$ for a constant unit vector $u$, so
    $c'(t) = s(t)\,u$ with $s := \lvert c'\rvert > 0$. Integrating from $0$ and using $c(0)=0$,
    \[c(t) = \sigma(t)\,u, \qquad \sigma(t) := \int_0^t s(\tau)\,d\tau,\]
    with $\sigma$ strictly increasing and $\sigma(0) = 0$. The endpoint condition $c(1) = x$
    reads $\sigma(1)\,u = x$, and $\sigma(1) = L(c)$ by definition of the length, so
    $L(c) = \lvert x\rvert$ and $u = x/\lvert x\rvert$ (note $x \neq 0$, since $\sigma(1) > 0$).
    The image is $\{\sigma u : 0 \leq \sigma \leq \lvert x\rvert\}$, the straight segment from
    $0$ to $x$.

    What the argument no longer pins down is the \emph{speed function} $s$, and with it the
    parametrization: any positive continuous $s$ with $\int_0^1 s = \lvert x\rvert$ produces a
    minimiser. That is unavoidable rather than a defect of the proof, because $L$ is invariant
    under reparametrization (\cref{rem:geodesic_constant_speed}), so it cannot distinguish two
    curves with the same image, and a criterion derived from $L$ alone cannot either. Assuming
    constant speed picks one representative out of that family, and $c'' \equiv 0$ is the
    statement for that representative.
\end{enumerate}
\end{exercisesolution}

\begin{exercisesolution}[Solution to \cref{ex:ai_cylinder_geodesics}]
% Generator: Claude Opus 5 (High)
\begin{enumerate}[label=\textbf{(\alph*)}]
  \item The first two coordinates satisfy $\cos^2(at)+\sin^2(at) = 1$ for every $t$, so
    $c(t) \in Z$. Differentiating,
    \[c'(t) = \big(-a\sin(at),\ a\cos(at),\ b\big), \qquad
      \lvert c'(t)\rvert = \sqrt{a^2\sin^2(at)+a^2\cos^2(at)+b^2} = \sqrt{a^2+b^2},\]
    which is constant and non-zero, since $a$ and $b$ are not both zero.
  \item Write $Z = F^{-1}\{0\}$ for $F(x,y,z) := x^2+y^2-1$. Then
    $\nabla F(x,y,z) = (2x,\,2y,\,0)$, which is non-zero on $Z$, so
    \cref{thm:regular_value_theorem} applies and \cref{def:normal_space} gives
    \[N_pZ = \operatorname{span}\{(p_1,\,p_2,\,0)\}, \qquad p \in Z.\]
    The cylinder's normal at a point is horizontal and radial: it forgets the height entirely.
    Differentiating $c$ a second time,
    \[c''(t) = \big(-a^2\cos(at),\ -a^2\sin(at),\ 0\big)
             = -a^2\big(c_1(t),\ c_2(t),\ 0\big),\]
    which is exactly a multiple of the spanning vector of $N_{c(t)}Z$. So
    $c''(t) \in N_{c(t)}Z$, i.e.\ $c''(t) \perp T_{c(t)}Z$, and the geodesic equation of
    \cref{prop:geodesic_equation_submanifold} holds for every $a$ and $b$.

    Note that the $z$-component of $c''$ vanishes identically, which is what makes this work.
    The vertical direction is tangential to $Z$ everywhere, so any vertical acceleration would
    have a tangential part and would violate the equation.
  \item With $a = 0$ the curve is $t \mapsto (1,0,bt)$, a vertical straight line ruling the
    cylinder. With $b = 0$ it is $t \mapsto (\cos(at),\sin(at),0)$, the horizontal unit circle.
    Both are limiting cases of the helix, and both are geodesics by \textbf{(b)}.
\end{enumerate}

\begin{remark}[Why the cylinder has \emph{only} these geodesics]
The three families found here are in fact all of them, and the quickest reason is that the
cylinder is \newterm{locally isometric} to the plane: rolling a sheet of paper into a cylinder
bends it without stretching, so it preserves the length of every curve drawn on it. Lengths are
all the geodesic equation ever sees, so geodesics must correspond to geodesics. Unrolling sends
helices to straight lines, vertical rulings to vertical lines, and horizontal circles to
horizontal lines, which is every straight line in the plane and nothing else.

The sphere admits no such trick, which is the same obstruction that makes flat world maps
distort areas or angles, and it is why \cref{ex:great_circles_geodesics} had to be checked by
computation instead.
\end{remark}
\end{exercisesolution}
